\section{线性响应理论}
\textrm{VASP}除了计算基态性质，还能计算系统对外界扰动的响应。如果扰动很小，响应位于线性响应区域，则适用线性响应理论。

\subsection{有限差分响应}
对\ch{SiC}，施加静态电场、应力和离子偏移，采用不同的有限差分方法，计算体系响应

根据微扰理论，受扰动体系的总能$E$可以在总能为$E_0$的未扰动体系的附近展开到两阶的电场\textcolor{red}{$\varepsilon$}、应力$\eta$和离子偏移$\mu$:
\begin{displaymath}
	\begin{aligned}
		E(\vec u, \varepsilon \varepsilon,\eta)=&E^0+\dfrac{\partial E}{\partial\mu_i}\mu_i+\dfrac{\partial E}{\partial\varepsilon_i}\varepsilon_i+\dfrac{\partial E}{\partial\eta_{ij}}\eta_{ij}\\
		+&\dfrac12\dfrac{\partial^2E}{\partial\mu_i\partial\mu_j}\mu_i\mu_j+\dfrac12\dfrac{\partial^2E}{\partial\varepsilon_i\partial\varepsilon_j}\varepsilon_i\varepsilon_j+\dfrac12\dfrac{\partial^2E}{\partial\eta_{ij}\partial\eta_{km}}\eta_{ij}\eta_{km}\\
		+&\dfrac{\partial^2E}{\partial\mu_i\partial\varepsilon_j}\mu_i\varepsilon_j+\dfrac{\partial^2E}{\partial\mu_i\partial\eta_{jk}\mu_i\eta_{jk}+\dfrac{\partial^2E}{\partial\varepsilon_i\partial\eta_{jk}}\varepsilon_i\eta_{jk}	
	\end{aligned}
\end{displaymath}
隐含了对下标的求和。注意，考虑到带电体系，对于外加电场，线性响应也会出现到另两个扰动项上的两阶上。

微扰的变分系数由其物理意义表示，即力$F$，电场极化$P$和应力$\sigma$写成
\begin{displaymath}
	\begin{aligned}
		F_i=&-\Omega_0\dfrac{\partial E}{\partial\mu_i}\bigg|_{\varepsilon,\eta}
		P_i=&-\dfrac{\partial E}{\partial\varepsilon_i}\bigg|_{\mu,\eta}=-\dfrac{fe}{(2\pi)^3}\sum_n\int_{\mathrm{BZ}}\mathrm{d}\vec k\langle \mu_{n\vec k}^{(\varepsilon_i)}|\mathrm{i}\nabla_{\vec k}|\mu_{n\vec k}^{(\varepsilon_i)}\rangle\\
		\sigma_{ij}=&\dfrac{\partial E}{\partial\eta_{ij}}\bigg|_{\mu,\varepsilon}
	\end{aligned}
\end{displaymath}
这里$\Omega_0$是未扰动体系的晶胞体积，$f=2$是未扰动计算的自旋多重度。$|\mu_{n,\vec k}^{(\varepsilon_i)}\rangle$是周期的场极化的\textrm{Kohn-Sham}轨道径向波函数。力常数矩阵$K$，冻结离子介电响应\textrm{(frozen ion dielectric susceptibility)}和离子钳位极化率\textrm{(ion-clamped polarizability)}\footnote{钳位\textrm{(clamped)}表示固定离子的边界条件。在\textrm{DFPT}计算中，这是计算介电张量的基本步骤，得到高频介电常数。}$\chi$，零外场弹性模量$C$表示为
\begin{displaymath}
	\begin{aligned}
		K_{ij}=&\Omega_0\dfrac{\partial^2E}{\partial\mu_i\partial\mu_j}\bigg|_{\varepsilon,\eta}\\
		\chi_{ij}=&-\dfrac{\partial^2E}{\partial\varepsilon_i\partial\varepsilon_j}\bigg|_{\mu,\eta}=\dfrac{\partial P_i}{\partial\varepsilon_j}\bigg|_{\mu,\eta}\\
		C_{ijkm}=&\Omega_0\dfrac{\partial^2E}{\partial\eta_{ij}\partial\eta_{km}\bigg|_{\mu,\varepsilon}
	\end{aligned}
\end{displaymath}
最终\textrm{Born}动力学有效电荷$Z^{\ast}$，(力响应)的内-应力张量$\Lambda$与离子钳位压电张量$e^{(0)}$可以写成
\begin{displaymath}
	\begin{aligned}
		Z_{ij}^{\ast}=&-\dfrac{\Omega_0}{e}\dfrac{\partial^2E}{\partial\mu_i\partial\varepsilon_j}a\bigg|_{\eta}=\dfrac{\Omega_0}{e}\dfrac{\partial F_i}{\partial\varepsilon_j}\bigg|_{\eta}\\
		\Lambda_{ijk}=&-\Omega_0\dfrac{\partial^2E}{\partial\mu_i\partial\varepsilon_{jk}}\bigg|_{\varepsilon}=\Omega_0\dfrac{\partial\sigma_{jk}}{\partial\mu_i}\bigg|_{\varepsilon}\\
		e_{ijk}^{(0)}=&-\dfrac{\partial^2E}{\partial\varepsilon_i\partial\eta_{jk}}\bigg|_{\mu}=\dfrac{\partial\sigma_{jk}}{\partial\varepsilon_i}\bigg|_{\mu}
\end{aligned}
\end{displaymath}
\textrm{VASP}中有两类不同层面的有限差分方法可供选择
\begin{itemize}
	\item 设置计算控制参数\texttt{IBRION}=5与6来完成离子偏移和应力计算
	\item 设置计算控制参数\texttt{LPEAD}来实施离散后扰动\textrm{(Perturbation Expression after discretization, PEAD)}方法，设置\texttt{LCALCEPS}来加外电场
\end{itemize}
注意:~该方法只能用于带隙体系，即不可以是导体或金属。通过离散型的\textrm{Sternheimer}方程计算与$\vec k$对应的\textrm{Kohn-Sham}轨道周期波函数的导数$\nabla_{\vec k}|\mu_{n\vec k}^{\varepsilon_i}\rangle$。这里的求和不包含空态。

\textrm{FCC}结构\ch{SiC}文件\textrm{POSCAR}

计算控制参数文件\textrm{INCAR}

电场极化计算控制参数文件\textrm{electric/INCAR}

电场极化-离子偏移计算控制参数文件\textrm{electric-ionic/INCAR}

$\vec k$空间网格剖分文件\textrm{KPOINTS}

能带计算高对称性$\vec k$点路径文件\textrm{KPOINTS}

第一个\textrm{POSCAR}文件表明计算采用默认的\textrm{PBE}泛函，计算控制参数\texttt{ISMEAR}和\texttt{SIGMA}指定\textrm{Gaussian}展宽方法处理\textrm{Fermi}面积分。计算中的能量截断设置得比较高以避免\textrm{Pulay}应力。当体系总能的能量差低于\texttt{EDIFF}则电子步弛豫收敛。

在电场极化计算\textrm{electric/INCAR}文件中，通过设置计算控制参数\texttt{LCALCPOL}=\textrm{T}，在\textrm{DFT}计算之初引入静态电场$\varepsilon_i$以计算电场极化。周期体系的电场极化的计算表达式为$-fe/(2\pi)^3$乘上对所有占据态的\textrm{Berry}相位$\int_{\mathrm{BZ}}\mathrm{d}\vec k\langle\mu_{n\vec k}^{(\varepsilon_i)}|\mathrm{i}\nabla|\mu_{n\vec k}^{(\varepsilon_i)}$的求和。

因为计算控制参数\texttt{LCALCEPS}=\textrm{T}，采用\textrm{PEAD}方法，因此对每个$\vec k$点的电场极化的\textrm{Kohn-Sham}轨道的径向部分有如下的有限差分来求导
\begin{displaymath}
	\dfrac{\partial|\mu_{n\vec k_j}^{(\varepsilon_i)}}{\partial k}=\dfrac1{2\Delta k}\sum_{m=1}^{N}\bigg[|\mu_{n\vec k_{j+1}}^{(\varepsilon_i)}\rangleS_{mn}^{-1}(\vec k_j,\vec k_{j+1})-|\mu_{n\vec k_{j-1}}^{(\varepsilon_i)}\rangleS_{mn}^{-1}(\vec k_j,\vec k_{j-1})\bigg]
\end{displaymath}
这里$S(\vec k_j,\vec k_{j+1})=\langle\mu_{n\vec k_j}^{(\varepsilon_i)}|\mu_{n\vec k_{j+1}}^{(\varepsilon_i)}\rangle$是重叠矩阵。

在电场极化-离子偏移计算控制参数\textrm{electric-ionic/INCAR}文件中，通过计算控制参数\texttt{IBROIN}额外增加离子偏移的有限差分计算。\texttt{NFREE}和\texttt{POTIM}分别表示离子步数和步长。\texttt{ISIF}=3允许晶胞体积和形状变化。这些都是弹性模量计算必须要的设置。

\textrm{KPOINTS}文件给出以$\Gamma$点为中心的均匀$\vec k$空间网格剖分，而\textrm{KPOINTS\_OPT}是用于能带计算的高对称性$\vec k$点路径布点。

应用\textrm{Bash}命令开始\textrm{VASP}计算\textrm{DFT}基态，通过\textrm{py4vasp}模块绘制能带图。

下一步，通过静态电场扰动基态体系。通过\textrm{Bash}命令，基于基态计算开始电场极化计算。

\textrm{VASP}计算会输出警示信息``The Electric field is too strong such that Zener tunneling can occur''。用户可以减少$\vec k$点或设置计算控制参数\texttt{EFIELD\_PEAD}来降低电场。下面的计算中，\texttt{EFIELD\_PEAD}会设置小一个数量级。

下一个计算中，加入可观的离子偏移量以改变晶胞体积和形状，以满足所有的静态线性响应性质计算。由基态\textrm{DFT}计算的\textrm{WAVECAR}出发，完成\textrm{Bash}名令执行的\textrm{VASP}计算。

然后通过\textrm{Py4vasp}模块提取线性响应的性质数据

将离子钳位极化率，弹性模量和压电张量与实验数据对比，这就是离子弛豫方法计算
\begin{displaymath}
	\begin{aligned}
		\chi_{ij}|_{\eta}=&\chi_{ij}|_{\mu,\eta}+\Omega_0^{(-1)}Z_{mi}(K^{-1})_{mn}Z_{nj}\\
		C_{ijkm}|_{\varepsilon}=&C_{ijkm}|_{\mu,\varepsilon}-\Omega_0^{(-1)}\Lambda_{nij}(K^{-1})_{no}\Lambda_{okm}\\
		e_{ijk}^{(0)}=&e_{ijk}^{(0)}|_{\mu}+\Omega_0^{(-1)}Z_{mi}(K^{-1})_{mn}\Lambda_{njk}
		\end{aligned}
\end{displaymath}
上述表达式隐含了对下标的求和。注意这里没有显式地给出极化率，因为\textrm{VASP}用极化率直接计算介电张量。%在实际计算中这样更方便，下一节会详细讨论。

在\textrm{OUTCAR}文件的末尾中可以找到相同的极化响应信息。

\subsection{密度泛函微扰理论与静态响应}
应用密度泛函微扰理论\textrm{(Density-Functional Perturbation Theory, DFPT)}计算忽略局域效应的长波极限下的\ch{AlP}的静态介电常数。

外加电场$\varepsilon_{\mathrm{ext},i}(\vec r^{\prime},\omega)$与介电常数作用，得到材料的内电场$\varepsilon_i(\vec r ,\omega)$，
\begin{displaymath}
	\varepsilon_i(\vec r,\omega)=\int\mathrm{d}\vec r^{\prime}\epsilon_{ij}^{-1}(\vec r,\vec r^{\prime},\omega)\varepsilon_{\mathrm{ext},j}(\vec r^{\prime,\omega})
\end{displaymath}
在动量空间内有
\begin{displaymath}
	\varepsilon_i(\vec q+\vec G,\omega)=\sum_{\vec G^{\prime}}\epsilon_{\vec G\vec G^{\prime}}^{(-1)}(\vec q,\omega)\varepsilon_{\mathrm{ext},j}(\vec q+\vec G^{\prime},\omega)
\end{displaymath}
如果忽略在$\vec r$和$\vec r^{\prime}$的细节区别，材料就是均匀的，外加电场的变化波长比原子尺度大很多。这就是长波近似\textrm{Long-wave length approximation}，此时$\epsilon_{ij}(\vec r-\vec r^{\prime,\omega})$仅与距离$\vec r-\vec r^{\prime}$有关。因此上式可以表示为
\begin{displaymath}
	\varepsilon_i(\vec q,\omega)=\epsilon_{(00),ij}^{-1}(\vec q,\omega)\varepsilon_{\mathrm{ext},j}(\vec q,\omega)
\end{displaymath}
换言之，当非局域场效应可以忽略时，$\vec G\neq\vec G^{\prime}$的非对角元为0。

离子钳位介电常数与离子钳位极化率密切关联，可有关系式
\begin{displaymath}
	\varepsilon_{ij}^{\infty}=\delat_{ij}+\dfrac{4\pi}{\varepsilon_0}\chi_{ij}
\end{displaymath}
这里$\delta_{ij}$时\textrm{Kronecker}符号，这里$\chi_{ij}$由上式确定，$\epsilon_0$时真空介电常数，$\epsilon^{\infty}$是对交变电流\textrm{(Alternating Current ,AC)}响应的介电函数，也可以是高于声子频率范围的光的电磁场频率。该频率对应于极限$\vec q\rightarrow0$，因此上标为$\infty$，在极化子研究中这是同行的标注方式。

\ch{AlP}的结构文件\textrm{POSCAR}

计算控制参数文件\textrm{IINCAR}

电场极化计算控制参数文件\textrm{electric-ionic/INCAR}

$\vec k$空间剖分网格文件\textrm{KPOINTS}

第一个计算的\textrm{INCAR}文件指定基本的\textrm{PBE}计算，而电场极化计算\textrm{electric-ionic/INCAR}文件中，采用\textrm{DFPT}计算，包括设置计算控制参数\texttt{LEPSILON}=\textrm{T}计算介电响应，设置\texttt{IBRION}=8计算离子偏移。

在\textrm{DFPT}理论中，径向周期\textrm{Kohn-Sham}波函数对电场$|\xi_{n\vec k}$的线性响应，可以由求解非微扰体系的线性\textrm{Sternheimer}方程确定
\begin{displaymath}
	[H(\vec k)-\epsilon_{n\vec k}S(\vec k)]|\xi_{n\vec k}=-\Delta H_{\mathrm{SCF}}(\veec k)|\mu_{n\vec k}\rangle-\hat{\mathbf q}\nabla_{\vec k}\mu_{n\vec k}
\end{displaymath}
这里$\Delta H_{\mathrm{SCF}}(\vec k)$是局域场效应，即因为\textrm{Kohn-Sham}波函数的变化，引起的微观周期体系的\textrm{Kohn-Sham}的\textrm{Hamiltonian}的改变。如果计算控制参数\texttt{LRPA}=\textrm{F}，那么$\Delta H_{\mathrm{SCF}}=0$。

在此基础上，基于\textrm{DFPT}，应用\textcolor{red}{两阶线性\textrm{Sternheimer}方程}
\begin{displaymath}
	[H(\vec k)-\epsilon_{n\vec k}S(\vec k)]|\nabla_{\vec k}\mu_{n\vec k}\rangle=\dfrac{\partial[H(\vec k)-\epsilon_{n\vec k}S(\vec k)]}{\partial\vec k}|\mu_{n\vec k}\rangle
\end{displaymath}
由此可以确定\textrm{Kohn-Sham}波函数对$\vec k$的导数。

最终静态宏观介电函数矩阵可以表示为
\begin{displaymath}
	\hat{\mathbf q}\cdot\epsilon^{\infty}\cdot\hat{\mathbf q}=1-\dfrac{f8\pi e^2}{\Omega_0}\sum_{n\vec k}w_{\vec k}\langle\hat{\mathbf q}\nabla_{\vec k}\mu_{n\vec k}|\xi_{n\vec k}\rangle
\end{displaymath}
\textcolor{red}{该表示式本质上与最初的表达式相同}，但是没有通过计算控制参数\texttt{LCALCEPS}选择有限差分方法来实现，而是通过\texttt{LEPSILON}选择了\textrm{DFPT}来计算得到响应函数。与\texttt{LCALCEPS}不同，\texttt{LEPSILON}是可以适用于导体和金属等周期体系的。两种计算方法都没有对未占据态的求和。目前\texttt{LEPSILON}还不能使用杂化泛函进行计算。%含有\textrm{Hartree-Fock}类

首先通过\textrm{Bash}命令完成基态的\textrm{VASP}计算。

在子目录\textrm{electric-ionic}下执行\textrm{Bash}命令完成介电响应计算。

在长波极限近似下，忽略局域场效应，宏观交变电场的介电响应函数$\epsilon^{\infty}$可以表示为:
\begin{displaymath}
	\epsilon_{ij}^{\infty}(\hat{\mathbf{q}},\omega)\approx\lim_{\mathbf{q}\rightarrow0}\epsilon_{00ij}(\mathbf{q},\omega)
\end{displaymath}
另一方面，如果局域场效应不能忽略，宏观交变电场的介电函数的倒数与$\mathbf{q}=0$极限的微观介电常数的关系为：
\begin{displaymath}
	\dfrac1{\epsilon_{ij}^{\infty}(\hat{\mathbf{q}},\omega)}\define\lim_{\mathbf{q}\rightarrow0}\epsilon_{00ij}^{-1}(\mathbf{q},\omega)
\end{displaymath}
应用\textrm{Py4vasp}模块给出静态介电响应函数，包括宏观静态介电张量、\textrm{Born}有效电荷和压电张量。

\subsection{频率相关的介电响应}
计算在独立粒子近似下\ch{NaCl}的频率相关介电函数。

完整的频率相关介电函数包括电子部分的和离子部分的贡献。在\textrm{DFT}计算中，这两部分分别计算后再加和
\begin{displaymath}
	\varepsilon(\omega)=\varepsilon_{\mathrm{elec}}(\omega)+\varepsilon_{\mathrm{ion}}(\omega)
\end{displaymath}
介电函数的虚部和实部满足\textrm{Kramer-Kronig}关系。从频率相关的介电函数出发，可以计算各类光学函数，如光反射率、光吸收率和光电导率等。

\textrm{FCC}结构的\ch{NaCl}文件\textrm{POSCAR}

计算控制文件\textrm{INCAR}

电子介电函数计算控制文件\textrm{electron/INCAR}

离子介电函数\textrm{DFPT}计算控制文件\textrm{ion-dfpt/INCAR}

离子介电函数有限差分计算控制文件\textrm{ion-finite-differences/INCAR}

$\vec k$空间网格剖分文件\textrm{KPOINTS}

周期体系的电子极化计算，是系数$-fe/(2\pi)^3$乘上\textrm{Berry}相位$\int_{\mathrm{BZ}}\mathrm{d}\vec k\langle\mu_{n\vec k}^{(\varepsilon_i)}|\mathrm{i}\nabla_{\vec k}|\mu_{n\vec k}^{(\varepsilon_i)}\rangle$对所有占据态的求和。电场极化的\textrm{Kohn-Sham}轨道波函数的径向部分对$\vec k$的导数$\nabla_{\vec k}|\mu_{n\vec k}^{(\varepsilon_i)}\rangle$，可以由两阶微扰理论计算：
\begin{displaymath}
	\nabla_{\vec k}|\mu_{n\vec k}^{(\varepsilon_i)}\rangle=\sum_{n^{\prime}\neq n}\dfrac{|\mu_{n^{\prime}\vec k}^{(\varepsilon_i)}\rangle\langle\mu_{n^{\prime}\vec k}^{(\varepsilon_i)}|\dfrac{\partial H(\vec k)-\epsilon_{n\vec k}S(\vec k)}{\partial\vec k}|\mu_{n\vec k}^{(\varepsilon_i)}}{\epsilon_{n\vec k}-\varepsilon_{n^{\prime}\vec k}}
\end{displaymath}
这里$H(\vec k)$是\textrm{Kohn-Sham}方程的径向波函数的\textrm{Hamiltonian}，能量是$\epsilon_{n\vec k}$，而$S(\vec k)$是重叠矩阵。计算控制参数\texttt{LOPTICS}=\textrm{T}时采用该表达式计算导数，这里的求和包含所有未占据态的求和。因此\texttt{NBANDS}必须包括足够多的未占据态。换言之，随着\texttt{NBANDS}数值的增加，极化将会趋于收敛。

这里采用的原子数据集选择适用于\textrm{GW}计算的，因为标准\textrm{DFT}计算中，未占据态的能带对基态性质，如体系总能的计算没有贡献。对于\textrm{GW}计算，原子数据集中，构造赝势时考虑了对激发态的贡献。因此当增加\texttt{NBANDS}时，应该对赝势适当优化，所以选择适用于\textrm{GW}计算的版本。

执行标准的\textrm{Bash}命令完成\textrm{VASP}计算，并通过\textrm{Py4vasp}计算\ch{NaCl}的带隙并绘制\textrm{DOS}曲线。

计算电子部分对频率相关的介电函数，将计算的\textrm{WAVECAR}复制到子目录\textrm{electron}，再次执行\textrm{Bash}命令，完成\textrm{VASP}计算。

通过\textrm{Bash}脚本\textrm{get\_dielectric\_function.sh}从\textrm{vasprun.xml}文件中提取介电函数数据，存入文件\textrm{dielectric\_function.dat}。通过\textrm{Py4vasp}绘制介电函数的曲线(包括实部与虚部)

当然，也可以直接通过\textrm{Py4vasp}绘制介电函数曲线。

离子部分对介电函数的贡献，可以选择采用\textrm{DFPT}方法或有限差分方法计算。分别在两个\textrm{INCAR}文件对应的目录下，通过\textrm{Bash}命令完成\textrm{VASP}计算，可以将两种计算的结果进行对比。

同样执行\textrm{Bash}脚本\textrm{get\_dielectric\_function.sh}，得到介电函数数据\textrm{dielectric\_function.dat}，然后通过\textrm{Py4vasp}模块绘制介电函数曲线。

有限差分方法计算耗时更长。

考虑到离子偏移，可以得到振动频率的计算表达式
\begin{displaymath}
	m_{\mathrm{ion}}\omega^2\mathbf{u}=\bigg(\dfrac{\partial^2E}{\partial u_i\partial u_j}\bigg)\mathbf{u}
\end{displaymath}
这里$\omega$是声子频率，$m_{\mathrm{ion}}$是离子质量，$\mathbf{u}$是离子偏移量，对应的\textrm{Hessian}矩阵与前面的相同。在\textrm{OUTCAR}文件中有声子频率的字符串提示：
\textsc{Eigenvectors and eigenvalues of the dynamical matrix}

声子频率与离子对介电函数的贡献的关系:~如果声子振动模具有偶极活性\textrm{(dipole-active)}\footnote{如果一个重复单元内，全部离子的电偶极矩对振动模有影响，那么该声子模就是偶极活性的。}，则该特征会出现在拼频率相关的介电函数中。这种情况出现在动力学矩阵的本征矢中。

除了三个偶极活性振动模，还有三个虚频率的声子模$(f/i)$。这里这些虚频振动模并不意味着结构是动力学不稳定的体系。采用不同的有限平面波基计算得到的体系总能略有不同，这些能量差别与晶胞和离子质心的相对位置有关。

对每个声子模，列出离子位置、质量权重和离子偏移。

两个矢量是平行的。但模量长度不同，这源于本征矢的质量权重导致的。这是离子质量出现在\textcolor{red}{方程(3.2)的后果。}离子质量选取\textrm{POTCAR}的计算控制参数\texttt{POMASS}。因此计算质量比
\begin{displaymath}
	\sqrt{\dfrac{m_{\mathrm{Na}}}{m_{\mathrm{Cl}}}}
\end{displaymath}
质量比与两个偏移矢量的比例相同。

给出的声子频率有各种单位。实际上，在这个算例中，离子和电子介电函数的频率轴也没有按照相同单位给出。

根据方程(3.2)，计算中频率的量纲单位是
\begin{displaymath}
	[\omega]=\sqrt{\dfrac1{[m_{\mathrm{ion}}]}\bigg[\dfrac{\partial^2E}{\partial u_i\partial u_j}\bigg]}
\end{displaymath}
\texttt{POMASS}按照原子单位\textrm{(atomic mass unit, a.m.u)}表示，有
\begin{displaymath}
	[m_{\mathrm{ion}}]=1.660599\cdot10^{-27}\mathrm{kg}
\end{displaymath}
能量单位是\textrm{eV}，离子偏移量的单位是\textrm{\AA}。所以频率单位可以写成\textrm{SI}
\begin{displaymath}
	1[\omega]=\sqrt{\dfrac{\mathrm{eV}/\AA^2}{\mathrm{a.m.u.}}}=9.822517\cdot10^{13}s^{-1}
\end{displaymath}
为了将计算的频率单位转换称\textrm{SI}单位，完成以下换算
\begin{displaymath}
\begin{aligned}
	\sqrt{\dfrac{\mathrm{eV}/\AA^2}{\mathrm{a.m.u.}}}=&\sqrt{\dfrac{1.602176487\cdot10^{-19}\mathrm{J}/(10^{-20}\mathrm{m}^2)}{1.660599\cdot10^{-27}\mathrm{kg}}}=\sqrt{\dfrac{16.02176487\mathrm{J/m^2}}{1.660599\cdot10^{-27}\mathrm{kg}}}\\
	=&9.822517\cdot10^{13}\mathrm{s}^{-1}
\end{aligned}
\end{displaymath}
对于另外的频率单位，如\textrm{eV}、\textrm{THz}和$\mathrm{cm}^{-1}$，也可以采用类似的换算。

为了得到\textrm{eV}，采用$E_{\mathrm{phonon}}=\hbar\omega$，因此有
\begin{displaymath}
	\begin{aligned}
	x[E_{\mathrm{phonon}}]=&1[\omega]\cdot\hbar=9.822517\cdot10^{13}\mathrm{s}^{-1}\times1.054572\cdot10^{-34}\mathrm{Js}=1.035855\codt10^{-20}\mathrm{J}\\
	=&0.064652976\mathrm{eV}
	\end{aligned}
\end{displaymath}
然后，角频率通过$\nu=\omeg/2\pi$换算
\begin{displaymath}
	x[\nv]=1[\omega]/2\pi=98.22517/(2\pi)\cdot10^{12}\mathrm{Hz}=15.6330214~\mathrm{THz}
\end{displaymath}
最终波长的倒数$\lambda^{-1}=\nv/c$，这里$c$是光速。由此得到
\begin{displaymath}
	x[\lambda^{-1}]=1[\nu]/c=\dfrac{15.6330214\cdot10^{12}\mathrm{s}^{-1}{29979245800\mathrm{cm/s}}}=521.46146389\mathrm{cm}^{-1}
\end{displaymath}
总之，频率单位的换算关系为
\begin{displaymath}
	\begin{algined}
		1\mathrm{THz}=&4.1357\mathrm{meV}=&33.356\mathrm{cm}^{-1}\\
		1\mathrm{meV}=&0.242\mathrm{THz}=&8.066\mathrm{cm}^{-1}\\
		1\mathrm{cm}^{-1}=&0.030\mathrm{THz}=&0.124\mathrm{meV}
	\end{algined}
\end{displaymath}
通过\textrm{Py4vasp}模块绘制\textrm{DFPT}方法计算离子对总的介电函数的贡献，频率单位是\textrm{meV}。

注意体系的总的最大频率来自离子的贡献。基于声子频率，离子的介电函数可以自动计算。具体到本算例，该值为$1.2\times\omega_{\mathrm{max}}$。

同样可以采用\textrm{Py4vasp}来绘制曲线。
